
\chapter{TRABALHOS RELACIONADOS}
\label{cap:correlatos}

O trabalho de \citeonline{xia_volearn_2022} teve como objetivo propor e validar o VoLearn, um sistema de aprendizagem de movimento cross-modal que permite a personalização de movimentos definidos pelo usuário, superando limitações de tutoriais convencionais baseados em movimentos predefinidos e feedback visual. O VoLearn utiliza tecnologias como a conversão de vídeos 2D em avatares 3D editáveis por meio do modelo DeepMotion, além de empregar o Unity3D para o desenvolvimento do ambiente virtual. O sistema integra um smartphone como sensor vestível, responsável por captar dados de rotação corporal e fornecer feedback auditivo ao usuário, composto por variações de tom, bips e instruções verbais. Sua aplicação prática abrange desde reabilitação física até treinamento esportivo e dança, permitindo que instrutores projetem movimentos a partir de vídeos e aprendizes pratiquem com orientação auditiva e visual. Os resultados dos estudos realizados indicaram que o feedback auditivo proporcionado pelo sistema foi eficaz na redução de erros de amplitude e tempo durante o aprendizado de movimentos, sem aumentar a carga cognitiva dos participantes. Entre as principais contribuições, destacam-se a flexibilidade para personalização dos movimentos, a acessibilidade proporcionada pelo uso de dispositivos cotidianos e a possibilidade de combinar diferentes movimentos para fins de coordenação e reabilitação. No entanto, o sistema apresenta limitações, como o suporte a apenas um sensor, dependência da qualidade do vídeo para reconstrução 3D, dificuldade na definição da complexidade dos movimentos e ausência de avaliação do aprendizado a longo prazo. Além disso, o ajuste de movimentos é restrito à amplitude tridimensional dos membros e a precisão do feedback depende da resposta individual ao áudio, não sendo validada a melhor posição do sensor para todos os usuários.

O estudo de \citeonline{emmanouil-two-2024} buscou determinar o número mínimo de repetições necessárias para garantir alta confiabilidade na medição do tempo de movimento em exercícios físicos, utilizando sensores inerciais. Foram utilizados sensores Xsens MTw Awinda, que combinam acelerômetros, giroscópios e magnetômetros, com os dados processados no MATLAB para análise temporal detalhada. Trinta adultos jovens e fisicamente ativos participaram do experimento, realizando oito tipos de exercícios fundamentais, enquanto um único sensor era posicionado no segmento corporal de maior amplitude de movimento. Os resultados mostraram que apenas duas repetições já são suficientes para alcançar excelente confiabilidade temporal, tanto dentro de uma mesma tentativa quanto entre diferentes tentativas, sendo quatro repetições o ideal para alta confiabilidade geral. As medições em segundos apresentaram maior confiabilidade relativa, enquanto as fases expressas em porcentagem do ciclo de movimento tiveram melhor confiabilidade absoluta. O estudo destaca que os sensores inerciais são ferramentas confiáveis, portáteis e econômicas para análise temporal de movimentos, otimizando o processo de coleta de dados em ambientes de treinamento e pesquisa. Entre as limitações, ressalta-se o foco em indivíduos saudáveis e ativos, o que pode restringir a generalização dos resultados para outras populações, além da dependência do tipo de exercício analisado. Ainda assim, a pesquisa contribui ao demonstrar que a avaliação temporal de movimentos com sensores inerciais pode ser realizada de forma prática e eficiente, simplificando protocolos de coleta e análise em contextos de fitness e treinamento.

O trabalho de \citeonline{pinto-aplicacao-2023} avaliou uma aplicação de Internet das Coisas (\gls{iot}) para identificar e monitorar movimentos em exercícios físicos com pesos livres, utilizando sensores inerciais e algoritmos de aprendizado de máquina. A solução proposta empregou sensores MPU-6050 acoplados a anilhas, microcontroladores ESP8266 com Wi-Fi, e comunicação via rede local, permitindo a coleta e transmissão dos dados para um servidor, onde foram processados e classificados por modelos de machine learning implementados em Python com scikit-learn. O sistema foi testado em diferentes exercícios com pesos livres, e os algoritmos de classificação (Árvore de Decisão, Perceptron Multicamadas e kNN) atingiram acurácias superiores a 90\% nos melhores cenários, especialmente quando utilizada maior dimensionalidade dos dados. A abordagem se destaca por ser de baixo custo, não invasiva e permitir treinamentos remotos, diferenciando-se de sistemas baseados em máquinas fixas ou dispositivos vestíveis. Entre as limitações, destaca-se que os resultados de alta acurácia foram obtidos apenas para dados de um único usuário, sendo necessária a coleta de dados de diferentes pessoas para garantir generalização. Além disso, a calibração dos sensores e a influência da gravidade representam desafios para a robustez do sistema. O estudo contribui para o avanço de soluções IoT aplicadas ao fitness, demonstrando a viabilidade de monitoramento e identificação automática de exercícios com pesos livres, e abre caminho para aplicações em saúde, como fisioterapia e acompanhamento remoto.

O estudo de \citeonline{dzaja-quantitative-nodate} propôs um procedimento abrangente para o monitoramento e avaliação quantitativa e qualitativa do movimento humano durante exercícios de força, utilizando dispositivos vestíveis equipados com sensores inerciais e magnéticos. Foram utilizados sensores Shimmer3 IMU posicionados no pulso, tórax e coxa, além de um módulo de ECG para monitoramento da frequência cardíaca, permitindo a coleta de dados detalhados em diferentes segmentos corporais. O sistema foi validado inicialmente em condições controladas com poucos sujeitos e, posteriormente, em um grupo maior e mais diverso, abrangendo diferentes níveis de experiência e faixas etárias. O método desenvolvido permite operar em três modos: básico (contagem de repetições), avançado (quantidade e qualidade dos movimentos, atuando como treinador virtual) e em tempo real (feedback imediato para sequências conhecidas). Foram implementados algoritmos de segmentação de repetições sem necessidade de conhecimento prévio do domínio, além de técnicas de classificação baseadas em machine learning, com destaque para o SVM, que atingiu acurácia superior a 99\% na identificação dos exercícios e das repetições. O sistema também propôs uma métrica para avaliação da qualidade do movimento, capaz de diferenciar sujeitos experientes de inexperientes e identificar repetições incorretas. O monitoramento da frequência cardíaca mostrou-se útil para avaliar o esforço e a qualidade do movimento. Entre as limitações, destaca-se o número reduzido de sujeitos na fase inicial, a necessidade de ampliar o conjunto de dados para incluir pessoas com lesões, e a recomendação de processos de anotação mais detalhados e avaliações mais abrangentes dos participantes. O estudo contribui ao oferecer um arcabouço flexível para avaliação automática e objetiva do desempenho em exercícios de força, com potencial para atuar como treinador virtual e apoiar profissionais de saúde e esporte no monitoramento e feedback em tempo real.

O estudo de \citeonline{chen-analysis-2024} investigou a biomecânica do agachamento com peso e desenvolveu um sistema de treinador físico baseado em inteligência artificial, utilizando deep learning e princípios da ciência do esporte para analisar a postura, calcular ângulos chave e identificar fases do movimento. O sistema empregou BlazePose e MediaPipe para estimação de pose, além de ferramentas como OpenCV, Matplotlib e Grafana para análise e visualização dos dados. Doze participantes experientes realizaram agachamentos sob diferentes condições de carga, permitindo a coleta de dados detalhados sobre ângulos articulares e padrões de movimento. Os resultados mostraram que o aumento do peso levou a maior inclinação do tronco e flexão profunda dos joelhos, com variações significativas nos ângulos ombro-quadril, joelho-quadril e joelho-tornozelo conforme a carga aumentava. O sistema de IA demonstrou eficácia ao fornecer feedback em tempo real, sugestões corretivas e relatórios personalizados, beneficiando os usuários com acompanhamento do progresso e visualização dos dados. Entre as limitações, destaca-se o tamanho reduzido da amostra e a necessidade de incluir participantes com diferentes níveis de experiência para ampliar a generalização dos resultados. O estudo contribui ao demonstrar o potencial da IA para aprimorar a técnica do agachamento, prevenir lesões e promover práticas de treinamento mais seguras e eficazes por meio de análise automatizada e feedback personalizado.

O estudo de \citeonline{gupta-yogahelp-2021} desenvolveu o sistema YogaHelp para auxiliar praticantes amadores a aprenderem a execução correta de exercícios de ioga sem a supervisão de um instrutor, reconhecendo os passos e avaliando o nível de correção da execução. O sistema utiliza quatro sensores de movimento (acelerômetro e giroscópio) acoplados aos membros, conectados a módulos NodeMCU (ESP32) e sensores IMU MPU9250, com coleta de dados via aplicativo Android e processamento em nuvem. Um modelo de deep learning baseado em redes neurais convolucionais é empregado para reconhecer os passos da Saudação ao Sol e avaliar a correção com base na velocidade de execução e nos ângulos dos membros, fornecendo feedback detalhado sobre velocidade e postura. O YogaHelp foi testado com instrutores e praticantes, demonstrando alta precisão no reconhecimento dos passos (acima de 96\% para instrutores e 93\% para praticantes) e melhora significativa na correção dos movimentos após quatro semanas de uso, com feedback eficaz para autoaperfeiçoamento. O sistema se mostrou superior a métodos existentes como SenseHAR e HuMAn, especialmente na precisão de reconhecimento e na redução da diferença de desempenho entre instrutores e praticantes. Entre as limitações, destaca-se o uso de dados de treino apenas de instrutores para o modelo, sugerindo a necessidade de incorporar aprendizagem incremental e explorar a integração com câmeras para futuras melhorias. O YogaHelp representa uma solução confiável para o ensino autônomo de ioga, especialmente relevante em contextos onde a supervisão presencial não é possível, validando a utilidade do feedback automatizado para o aprendizado de posturas corretas.

O estudo de \citeonline{balkhi-multipurpose-nodate} apresentou uma plataforma vestível multifuncional baseada em sensores, desenvolvida para detecção automática de peso, reconhecimento do tipo de atividade e contagem de repetições em treinamentos de força. A solução consiste em uma luva inteligente equipada com sensores de força, acelerômetro, IMU e comunicação Bluetooth, integrando Arduino Nano e um aplicativo móvel para coleta e visualização dos dados. O sistema foi testado com participantes realizando diferentes exercícios e pesos, demonstrando alta acurácia na detecção de peso (SVM: 99,8\%), reconhecimento de atividades (Random Forest: 99,8\%) e contagem de repetições (96\%). A plataforma se destaca por consolidar múltiplas análises em um único dispositivo, reduzindo o desconforto do usuário e fornecendo feedback automático sobre o progresso, prevenindo excesso de esforço e lesões. Entre as limitações, estão o desafio de coletar dados para toda a variedade de atividades esportivas, a necessidade de investigar a detecção de pesos mais pesados e diferentes formatos de halteres, o desvio dos sensores de força com uso prolongado e o fato de a luva protótipo ser apenas para a mão direita. O estudo contribui ao oferecer uma solução prática e eficaz para monitoramento em tempo real de treinos de força, com potencial para prevenir lesões e otimizar o desempenho em ambientes de academia.

O estudo de \citeonline{hwang-estimation-2023} desenvolveu e validou um modelo para estimar o one-repetition maximum (1-RM) em exercícios com faixas de resistência, além de uma rede neural convolucional (CNN) para identificar o tipo de exercício e contar repetições, utilizando dados de câmera RGB e sensores IMU. Trinta participantes realizaram cinco exercícios de membros superiores, com coleta de dados por vídeo, sensores inerciais e análise por OpenPose, além de testes com halteres e faixas de diferentes resistências. O modelo CNN que combinou dados do IMU e das articulações corporais atingiu precisão de 98,83\% na classificação dos exercícios, enquanto o algoritmo de contagem de repetições apresentou precisão variável conforme o exercício, sendo mais preciso para supino, desenvolvimento de ombros e remada sentada. As equações de regressão para estimativa do 1-RM mostraram alta correlação entre a força da faixa e o 1-RM, com R² ajustado acima de 0,84 para todos os exercícios. Entre as limitações, destacam-se a restrição a cinco exercícios de membros superiores, a simplicidade do modelo de CNN para capturar variações temporais e imprecisões na contagem de repetições em movimentos não lineares ou com poucas repetições. O estudo contribui ao propor uma solução portátil e de baixo custo para avaliação da força muscular e monitoramento de exercícios com faixas de resistência, relevante para o mercado crescente de treinamento sem contato, e destaca a necessidade de pesquisas futuras para ampliar a aplicabilidade e aprimorar a avaliação da qualidade e intensidade dos exercícios.

A análise dos trabalhos correlatos permite identificar diferentes estratégias utilizadas para o monitoramento de exercícios. \citeonline{xia_volearn_2022}, com o sistema Volearn, destacam-se pela proposta de fornecer feedback auditivo em tempo real, apontando a importância de não sobrecarregar o usuário com informações excessivas. \citeonline{emmanouil-two-2024} investigaram a confiabilidade de medições em exercícios, mostrando o número mínimo de repetições necessárias para resultados consistentes, o que contribui para a definição de protocolos de coleta de dados. \citeonline{pinto-aplicacao-2023} avançaram ao propor a instrumentação direta de anilhas e a integração com IoT, reduzindo o desconforto dos sensores vestíveis e trazendo inspiração direta para este trabalho.  

Outros estudos exploraram diferentes perspectivas. \citeonline{dzaja-quantitative-nodate} reforçaram o papel do feedback em tempo real e apresentaram análises de algoritmos de aprendizado de máquina aplicados ao movimento humano. \citeonline{chen-analysis-2024} introduziram o uso de visão computacional e pose estimation, permitindo a geração de relatórios detalhados após o treino, enquanto \citeonline{hwang-estimation-2023} ampliaram essa abordagem com uso de imagens RGB e OpenPose para análises comparativas. \citeonline{gupta-yogahelp-2021}, por sua vez, trouxeram a ideia de aplicativos específicos para treinadores rotularem execuções corretas, proposta que inspirou diretamente a adaptação aqui apresentada, na forma de um modo treinador dentro do aplicativo. Por fim, \citeonline{balkhi-multipurpose-nodate} mostraram que o uso de uma luva com sensores IMU atingiu alta acurácia na detecção de atividades e contagem de repetições, demonstrando que é possível realizar reconhecimento mesmo com informações limitadas, mas evidenciando também o incômodo associado a soluções vestíveis.  

Dessa forma, observa-se que cada um desses trabalhos contribui de forma complementar para o avanço da área, mas ainda não há uma solução que una portabilidade, baixo custo, ausência de desconforto, feedback em tempo real e personalização do aprendizado. É nesse espaço que o presente trabalho se insere, ao propor halteres instrumentados com IMUs e um aplicativo com dois modos de operação, capazes de oferecer feedback automático ao usuário e, simultaneamente, permitir que treinadores alimentem o algoritmo com informações específicas sobre a execução.

Os Quadros \ref{tab:correlatos1}, \ref{tab:correlatos2} e \ref{tab:correlatos3} apresentam uma síntese comparativa dos principais trabalhos relacionados, destacando as características técnicas, as funcionalidades e as metodologias empregadas. Observa-se que a maioria das propostas utiliza sensores vestíveis, em alguns casos associados a outros dispositivos, como câmeras, o que pode aumentar o custo ou o desconforto do usuário. Quanto às funcionalidades, há soluções que reconhecem exercícios e contam repetições, mas apenas parte delas realiza a avaliação da qualidade do movimento em tempo real, aspecto considerado para prevenção de lesões e para a eficácia do treinamento. Além disso, embora alguns estudos empreguem técnicas de aprendizado profundo e visão computacional, a dependência de câmeras limita a portabilidade e a aplicabilidade em ambientes não controlados.  

Com isso, o trabalho aqui proposto diferencia-se por integrar sensores inerciais diretamente aos halteres, eliminando a necessidade de vestíveis adicionais ou câmeras, ao mesmo tempo em que prioriza a avaliação qualitativa da execução com suporte de algoritmos de aprendizado profundo. Essa abordagem busca unir praticidade, baixo custo, portabilidade e personalização do feedback, atendendo tanto ao praticante, por meio de correções automáticas em tempo real, quanto ao profissional, que pode supervisionar e enriquecer a base de dados do sistema.

\begin{quadro}[H]
\centering
\caption{Comparação dos trabalhos quanto ao uso de equipamentos e sensores}
\label{tab:correlatos1}
\renewcommand{\arraystretch}{1.3}
\begin{tabular}{|l|C{2cm}|C{2cm}|C{2cm}|}
\hline
\textbf{Trabalho} & \textbf{É vestível?} & \textbf{Sensores além de IMU?} & \textbf{Feedback em tempo real?} \\
\hline
\cite{balkhi-multipurpose-nodate}            & Sim & Sim & Não \\\hline
\cite{chen-analysis-2024}       & Não & Não & Sim \\\hline
\cite{dzaja-quantitative-nodate}         & Sim & Sim & Sim \\\hline
\cite{emmanouil-two-2024}  & Sim & Não & Não \\\hline
\cite{gupta-yogahelp-2021}    & Sim & Não & Sim \\\hline
\cite{hwang-estimation-2023}      & Sim & Sim & Não \\\hline
\cite{pinto-aplicacao-2023}             & Não & Não & Não \\\hline
\cite{xia_volearn_2022}     & Sim & Não & Sim \\\hline
\textbf{Este trabalho}        & Não & Não & Sim \\
\hline
\end{tabular}
\end{quadro}

\begin{quadro}[H]
\centering
\caption{Comparação dos trabalhos quanto ao reconhecimento e análise dos exercícios}
\label{tab:correlatos2}
\renewcommand{\arraystretch}{1.3}
\begin{tabular}{|l|C{2cm}|C{2cm}|C{2cm}|}
\hline
\textbf{Trabalho} & \textbf{Reconhecimento de exercícios?} & \textbf{Contagem de repetições?} & \textbf{Avalia qualidade do movimento?} \\
\hline
\cite{balkhi-multipurpose-nodate}            & Sim & Sim & Não \\\hline
\cite{chen-analysis-2024}       & Sim & Sim & Sim \\\hline
\cite{dzaja-quantitative-nodate}         & Sim & Sim & Sim \\\hline
\cite{emmanouil-two-2024}  & Não & Não & Não \\\hline
\cite{gupta-yogahelp-2021}    & Sim & Não & Sim \\\hline
\cite{hwang-estimation-2023}      & Sim & Sim & Não \\\hline
\cite{pinto-aplicacao-2023}             & Sim & Não & Não \\\hline
\cite{xia_volearn_2022}         & Sim & Não & Sim \\\hline
\textbf{Este trabalho}        & - & Sim & Sim \\
\hline
\end{tabular}
\end{quadro}

\begin{quadro}[H]
\centering
\caption{Comparação dos trabalhos quanto a testes, uso de câmeras e técnicas de aprendizado profundo}
\label{tab:correlatos3}
\renewcommand{\arraystretch}{1.3}
\begin{tabular}{|l|C{2cm}|C{2cm}|C{2cm}|}
\hline
\textbf{Trabalho} & \textbf{Testado em exercícios de força?} & \textbf{Utiliza câmera?} & \textbf{Aplica Deep Learning?} \\
\hline
\cite{balkhi-multipurpose-nodate}            & Sim & Não & Não \\\hline
\cite{chen-analysis-2024}       & Sim & Sim & Sim \\\hline
\cite{dzaja-quantitative-nodate}         & Sim & Não & Não \\\hline
\cite{emmanouil-two-2024}  & Sim & Não & Não \\\hline
\cite{gupta-yogahelp-2021}    & Sim & Não & Sim \\\hline
\cite{hwang-estimation-2023}      & Sim & Sim & Sim \\\hline
\cite{pinto-aplicacao-2023}             & Sim & Não & Não \\\hline
\cite{xia_volearn_2022}         & Sim & Sim & Sim \\\hline
\textbf{Este trabalho}        & Sim & Não & Sim \\
\hline
\end{tabular}
\end{quadro}

